% Options for packages loaded elsewhere
\PassOptionsToPackage{unicode}{hyperref}
\PassOptionsToPackage{hyphens}{url}
\documentclass[
  ignorenonframetext,
]{beamer}
\newif\ifbibliography
\usepackage{pgfpages}
\setbeamertemplate{caption}[numbered]
\setbeamertemplate{caption label separator}{: }
\setbeamercolor{caption name}{fg=normal text.fg}
\beamertemplatenavigationsymbolsempty
% remove section numbering
\setbeamertemplate{part page}{
  \centering
  \begin{beamercolorbox}[sep=16pt,center]{part title}
    \usebeamerfont{part title}\insertpart\par
  \end{beamercolorbox}
}
\setbeamertemplate{section page}{
  \centering
  \begin{beamercolorbox}[sep=12pt,center]{section title}
    \usebeamerfont{section title}\insertsection\par
  \end{beamercolorbox}
}
\setbeamertemplate{subsection page}{
  \centering
  \begin{beamercolorbox}[sep=8pt,center]{subsection title}
    \usebeamerfont{subsection title}\insertsubsection\par
  \end{beamercolorbox}
}
% Prevent slide breaks in the middle of a paragraph
\widowpenalties 1 10000
\raggedbottom
\AtBeginPart{
  \frame{\partpage}
}
\AtBeginSection{
  \ifbibliography
  \else
    \frame{\sectionpage}
  \fi
}
\AtBeginSubsection{
  \frame{\subsectionpage}
}
\usepackage{iftex}
\ifPDFTeX
  \usepackage[T1]{fontenc}
  \usepackage[utf8]{inputenc}
  \usepackage{textcomp} % provide euro and other symbols
\else % if luatex or xetex
  \usepackage{unicode-math} % this also loads fontspec
  \defaultfontfeatures{Scale=MatchLowercase}
  \defaultfontfeatures[\rmfamily]{Ligatures=TeX,Scale=1}
\fi
\usepackage{lmodern}
\ifPDFTeX\else
  % xetex/luatex font selection
\fi
% Use upquote if available, for straight quotes in verbatim environments
\IfFileExists{upquote.sty}{\usepackage{upquote}}{}
\IfFileExists{microtype.sty}{% use microtype if available
  \usepackage[]{microtype}
  \UseMicrotypeSet[protrusion]{basicmath} % disable protrusion for tt fonts
}{}
\makeatletter
\@ifundefined{KOMAClassName}{% if non-KOMA class
  \IfFileExists{parskip.sty}{%
    \usepackage{parskip}
  }{% else
    \setlength{\parindent}{0pt}
    \setlength{\parskip}{6pt plus 2pt minus 1pt}}
}{% if KOMA class
  \KOMAoptions{parskip=half}}
\makeatother
\usepackage{longtable,booktabs,array}
\newcounter{none} % for unnumbered tables
\usepackage{calc} % for calculating minipage widths
\usepackage{caption}
% Make caption package work with longtable
\makeatletter
\def\fnum@table{\tablename~\thetable}
\makeatother
\usepackage{graphicx}
\makeatletter
\newsavebox\pandoc@box
\newcommand*\pandocbounded[1]{% scales image to fit in text height/width
  \sbox\pandoc@box{#1}%
  \Gscale@div\@tempa{\textheight}{\dimexpr\ht\pandoc@box+\dp\pandoc@box\relax}%
  \Gscale@div\@tempb{\linewidth}{\wd\pandoc@box}%
  \ifdim\@tempb\p@<\@tempa\p@\let\@tempa\@tempb\fi% select the smaller of both
  \ifdim\@tempa\p@<\p@\scalebox{\@tempa}{\usebox\pandoc@box}%
  \else\usebox{\pandoc@box}%
  \fi%
}
% Set default figure placement to htbp
\def\fps@figure{htbp}
\makeatother
\setlength{\emergencystretch}{3em} % prevent overfull lines
\providecommand{\tightlist}{%
  \setlength{\itemsep}{0pt}\setlength{\parskip}{0pt}}
\usepackage{bookmark}
\IfFileExists{xurl.sty}{\usepackage{xurl}}{} % add URL line breaks if available
\urlstyle{same}
\hypersetup{
  hidelinks,
  pdfcreator={LaTeX via pandoc}}

\author{\texorpdfstring{}{}}
\date{}

\begin{document}

\begin{frame}
\newpage
\end{frame}

\begin{frame}[fragile]{Experimental Validation}
\protect\phantomsection\label{sec:results}
This section demonstrates the practical viability of CIF's formal
mechanisms through empirical evaluation across production multiagent
architectures. We present experimental setup (\cref{sec:exp-setup}) and
key findings (\cref{sec:key-findings}). Detailed statistical analysis,
ablation studies, and scalability metrics are provided in
\cref{sec:extended-results}.

\begin{block}{Experimental Setup}
\protect\phantomsection\label{sec:exp-setup}
\begin{block}{Target Architectures}
\protect\phantomsection\label{target-architectures}
We evaluated CIF across six production multiagent systems representing
diverse architectural patterns:

\textbf{Table: Multiagent system architectures evaluated.}
\{\#tab:target-architectures\}

{\def\LTcaptype{none} % do not increment counter
\begin{longtable}[]{@{}lll@{}}
\toprule\noalign{}
System & Architecture & Communication \\
\midrule\noalign{}
\endhead
Claude Code & Hierarchical (\(1 + n\)) & Task delegation \\
AutoGPT & Autonomous + plugins & Tool-based \\
CrewAI & Role-based (3--10) & Sequential/parallel \\
LangGraph & Graph-based & State machine \\
MetaGPT & SOP-driven (5--8) & Document passing \\
Camel & Debate (\(2+\)) & Adversarial \\
\bottomrule\noalign{}
\end{longtable}
}
\end{block}

\begin{block}{Attack Corpus}
\protect\phantomsection\label{attack-corpus}
We assembled a corpus of 950 cognitive attacks across four categories:
prompt injection (500), trust exploitation (200), belief manipulation
(150), and coordination attacks (100). Sources include published
jailbreak datasets \cite{perez2023hackaprompt}, custom adversarial
prompts, red team exercises (8 researchers, 4-week engagement), and
synthetic generation via adversarial models. The corpus was partitioned
into training (70\%, \(n=665\)), test (20\%, \(n=190\)), and validation
(10\%, \(n=95\)) splits with stratification across attack categories and
difficulty levels. The training set was used for threshold calibration
and fusion operator training; the test set for primary performance
evaluation; and the held-out validation set for generalization
assessment and sensitivity analysis.
\end{block}

\begin{block}{Evaluation Methodology}
\protect\phantomsection\label{sec:eval-methodology}
\textbf{Simulation-Based Analysis.} We evaluate CIF through parametric,
architecture-aware simulation rather than live deployment against
production systems. Each target architecture is modeled as a topology
adapter that captures three structural properties: (1) communication
pattern (hierarchical, peer-to-peer, role-based, graph-based), (2) trust
structure (centralized authority, distributed reputation, role-based
permissions), and (3) attack surface characteristics (entry points,
propagation paths, state exposure). For each attack sample, the
evaluation framework computes a detection score from calibrated base
rates (indexed by attack difficulty: easy, medium, hard), modulated by
architecture-specific attack-surface multipliers that reflect the
topology's structural exposure, with Gaussian noise (\(\sigma = 0.05\))
added for stochastic variation. The resulting score is thresholded
(\(\tau = 0.5\)) to produce a binary detection outcome.

This parametric approach enables controlled comparison across
architectures and attack types at scale (950 \(\times\) 6 = 5,700
evaluation instances), systematic sensitivity analysis across parameter
configurations, and ablation studies that would be prohibitively
expensive against live systems. The base detection rates and
architecture multipliers were calibrated against published benchmarks
for prompt injection detection
\cite{greshake2023indirect, liu2023prompt} and the authors' experience
deploying cognitive firewalls in production multiagent systems.
Crucially, the reported detection rates characterize the
\emph{framework's design-level detection properties under calibrated
conditions}, not the output of running attack text through the
implemented defense modules in real time.

\textbf{Relationship to implemented modules.} The CIF defense
modules---Cognitive Firewall, Belief Sandbox, Tripwire Monitor, Trust
Calculus, Byzantine Consensus, Provenance, Drift Detection, and
Invariant Checker---are fully implemented and independently tested
(1,557 unit and integration tests at 100\% pass rate). The evaluation
framework's \texttt{ExperimentRunner} accepts an optional defense
pipeline; when a real pipeline is provided, it routes each attack
sample's text content through the actual defense modules (e.g.,
\texttt{EnhancedCognitiveFirewall.classify\textbackslash{}\_detailed()},
\texttt{CognitiveTripwire.check()}) instead of the parametric model.
This design enables direct recomposition: replacing simulation with
pipeline-driven evaluation requires only passing the assembled defense
pipeline to the runner, with no changes to the evaluation harness,
attack corpus, or analysis scripts.

\textbf{Limitations of the simulation approach.} Because the evaluation
uses calibrated parametric simulation rather than pipeline-driven
detection, the reported rates reflect CIF's design-level detection
properties---how the defense layers \emph{should} perform given
calibrated difficulty and architecture characteristics---rather than
measured outcomes from processing attack text through the real modules.
Native defenses built into each framework (e.g., Claude Code's
permission gating, AutoGPT's safety constraints) are not captured in the
baseline, which assumes no CIF components are active. Results should
therefore be interpreted as characterizing CIF's intrinsic detection
architecture rather than marginal improvement over existing framework
protections. The immediate next step for empirical validation is to run
the full 950-attack corpus through the real defense pipeline and compare
pipeline-driven detection rates against the calibrated simulation
baseline.

\textbf{Operational Definition of Detection.} An attack is classified as
\texttt{detected\textquotesingle{}\textquotesingle{}\ when\ the\ CIF\ pipeline\textquotesingle{}s\ aggregate\ confidence\ score\ exceeds\ the\ configured\ threshold\ (\$\textbackslash{}tau\ =\ 0.5\$\ by\ default).\ The\ confidence\ score\ combines\ firewall\ pattern-matching\ scores,\ sandbox\ quarantine\ signals,\ tripwire\ alerts,\ and\ trust\ calculus\ violations\ via\ the\ learned\ fusion\ operator\ (Section\ 2).\ Ground\ truth\ labels\ were\ assigned\ by\ two\ independent\ annotators\ (Cohen\textquotesingle{}s\ \$\textbackslash{}kappa\ =\ 0.84\$,\ indicating}almost
perfect'\,' agreement per \cite{landis1977measurement}) with
disagreements resolved by a third reviewer.

\textbf{Reproducibility.} All experiments use a fixed random seed (42)
for deterministic reproduction. The complete evaluation framework,
including architecture adapters, attack corpus (sanitized subset), and
analysis scripts, is available in the supplementary repository.
Multi-seed stability analysis across 30 seeds is reported in
\cref{sec:extended-results}. All experiments are fully deterministic
when executed with the default seed configuration.

\textbf{Runtime and Resource Requirements.} The full experiment suite
completes in approximately 15 minutes on the reference hardware
specified below. Peak memory usage reaches approximately 2GB during
Byzantine consensus tests, which require maintaining state for all agent
interactions. Individual defense mechanism tests (firewall-only,
sandbox-only, tripwires-only) complete in under 5 minutes each.

\textbf{Software Environment.} Python 3.12, NumPy 1.26, SciPy 1.12,
scikit-learn 1.4, matplotlib 3.8. The evaluation framework has been
tested on Python 3.10, 3.11, and 3.12 with consistent results across all
versions. All experiments executed on a single workstation (Apple M3
Max, 128GB RAM, macOS 15).
\end{block}
\end{block}

\begin{block}{Key Findings}
\protect\phantomsection\label{sec:key-findings}
\begin{block}{Finding 1: Layered Defense Significantly Outperforms
Single Mechanisms}
\protect\phantomsection\label{finding-1-layered-defense-significantly-outperforms-single-mechanisms}
The central empirical finding validates CIF's layered approach. No
single defense mechanism achieves acceptable protection, but their
composition yields substantial improvement.
\Cref{fig:detection-performance} presents detection rates across defense
configurations and attack categories.

\begin{figure}
\centering
\includegraphics[width=0.95\linewidth,height=\textheight,keepaspectratio,alt={Detection Performance Comparison. Bar chart comparing detection rates across defense configurations (Baseline, Firewall-only, Sandbox-only, Tripwires-only, Full CIF) for each attack category (Prompt Injection, Trust Exploitation, Belief Manipulation, Coordination). Error bars show 95\% bootstrap confidence intervals (n=1\{,\}000 resamples). Full CIF consistently achieves \textgreater90\textbackslash\% detection across all categories, while individual mechanisms show significant gaps---validating the defense composition algebra (Part 1, Theorems 3.1-3.2). The Full CIF theoretical rate (1 - \textbackslash prod(1-r\_i) \textbackslash approx 0.99) via  exceeds the empirical 0.94, indicating room for implementation-level optimization. Detection data generated by the CIF evaluation pipeline ().}]{../figures/detection_performance.pdf}
\caption{Detection Performance Comparison. Bar chart comparing detection
rates across defense configurations (Baseline, Firewall-only,
Sandbox-only, Tripwires-only, Full CIF) for each attack category (Prompt
Injection, Trust Exploitation, Belief Manipulation, Coordination). Error
bars show 95\% bootstrap confidence intervals (\(n=1{,}000\) resamples).
Full CIF consistently achieves \(>90\%\) detection across all
categories, while individual mechanisms show significant
gaps---validating the defense composition algebra (Part 1, Theorems
3.1-3.2). The Full CIF theoretical rate
(\(1 - \prod(1-r_i) \approx 0.99\)) via
\texttt{compute\_series\_detection\_rate()} exceeds the empirical 0.94,
indicating room for implementation-level optimization. Detection data
generated by the CIF evaluation pipeline
(\texttt{output/data/detection\_data.json}).}\label{fig:detection-performance}
\end{figure}

\textbf{Table: Detection performance by defense configuration.}
\{\#tab:detection-performance\}

{\def\LTcaptype{none} % do not increment counter
\begin{longtable}[]{@{}
  >{\raggedright\arraybackslash}p{(\linewidth - 12\tabcolsep) * \real{0.1429}}
  >{\raggedright\arraybackslash}p{(\linewidth - 12\tabcolsep) * \real{0.1429}}
  >{\raggedright\arraybackslash}p{(\linewidth - 12\tabcolsep) * \real{0.1429}}
  >{\raggedright\arraybackslash}p{(\linewidth - 12\tabcolsep) * \real{0.1429}}
  >{\raggedright\arraybackslash}p{(\linewidth - 12\tabcolsep) * \real{0.1429}}
  >{\raggedright\arraybackslash}p{(\linewidth - 12\tabcolsep) * \real{0.1429}}
  >{\raggedright\arraybackslash}p{(\linewidth - 12\tabcolsep) * \real{0.1429}}@{}}
\toprule\noalign{}
\begin{minipage}[b]{\linewidth}\raggedright
Defense
\end{minipage} & \begin{minipage}[b]{\linewidth}\raggedright
Prompt Inj.
\end{minipage} & \begin{minipage}[b]{\linewidth}\raggedright
Trust Expl.
\end{minipage} & \begin{minipage}[b]{\linewidth}\raggedright
Belief Manip.
\end{minipage} & \begin{minipage}[b]{\linewidth}\raggedright
Coord.
\end{minipage} & \begin{minipage}[b]{\linewidth}\raggedright
Overall
\end{minipage} & \begin{minipage}[b]{\linewidth}\raggedright
Key Limitation
\end{minipage} \\
\midrule\noalign{}
\endhead
Baseline (no CIF)\(^\dagger\) & 0.00 & 0.00 & 0.00 & 0.00 & 0.00 & No
defense active \\
Firewall only & 0.85 & 0.62 & 0.70 & 0.55 & 0.74 & Misses coordination
attacks \\
Sandbox only & 0.69 & 0.72 & 0.78 & 0.62 & 0.71 & Limited to unverified
sources \\
Tripwires only & 0.82 & 0.87 & 0.86 & 0.79 & 0.83 & Requires canary
placement \\
\textbf{Full CIF} & \textbf{0.96} & \textbf{0.92} & \textbf{0.94} &
\textbf{0.90} & \textbf{0.94} {[}0.92, 0.96{]} & Acceptable latency
overhead \\
\bottomrule\noalign{}
\end{longtable}
}

\(^\dagger\)\textit{Baseline reflects no CIF components active, not the absence of all defenses---production frameworks include native safety features that would provide non-zero baseline protection. The baseline condition represents agent performance under each attack category with no defensive mechanisms active, establishing the undefended success rate against which CIF improvements are measured.}

The gap between firewall-only and full CIF is most pronounced for
coordination and temporal attacks, which require multi-component
detection. This validates the defense composition algebra (Section 4
(Defense Composition, Part 1)): defenses targeting orthogonal attack
surfaces compose multiplicatively.
\end{block}

\begin{block}{Finding 2: Trust Calculus Prevents Amplification Attacks}
\protect\phantomsection\label{finding-2-trust-calculus-prevents-amplification-attacks}
\begin{figure}
\centering
\includegraphics[width=0.9\linewidth,height=\textheight,keepaspectratio,alt={ROC Curves by Attack Category. Receiver Operating Characteristic curves showing the tradeoff between True Positive Rate (sensitivity) and False Positive Rate (1-specificity) for CIF detection across four attack categories. AUC values computed via  (trapezoidal integration): Prompt Injection = 0.97, Trust Exploitation = 0.94, Belief Manipulation = 0.93, Coordination Attacks = 0.92. The operating point (diamond marker) is selected to maximize Youden's J statistic (J = \textbackslash text\{TPR\} - \textbackslash text\{FPR\}), balancing sensitivity and specificity. Confidence bands (shaded) computed via n=1\{,\}000 bootstrap resamples; the wider band for Coordination Attacks reflects the smaller sample size (n=100 vs.~n=500 for injection).}]{../figures/roc_curves.pdf}
\caption{ROC Curves by Attack Category. Receiver Operating
Characteristic curves showing the tradeoff between True Positive Rate
(sensitivity) and False Positive Rate (1-specificity) for CIF detection
across four attack categories. AUC values computed via
\texttt{compute\_auc\_from\_points()} (trapezoidal integration): Prompt
Injection = 0.97, Trust Exploitation = 0.94, Belief Manipulation = 0.93,
Coordination Attacks = 0.92. The operating point (diamond marker) is
selected to maximize Youden's \(J\) statistic
(\(J = \text{TPR} - \text{FPR}\)), balancing sensitivity and
specificity. Confidence bands (shaded) computed via \(n=1{,}000\)
bootstrap resamples; the wider band for Coordination Attacks reflects
the smaller sample size (\(n=100\) vs.~\(n=500\) for
injection).}\label{fig:roc-curves}
\end{figure}

The ROC analysis (\cref{fig:roc-curves}) demonstrates strong
discrimination across all attack categories. Across all tested
architectures, the bounded trust decay (\(\delta^d\)) successfully
prevented trust laundering and amplification attempts. In adversarial
scenarios where attackers attempted to relay high-impact content through
multiple trusted intermediaries, the exponential decay ensured that
delegated trust remained below action thresholds.

Critically, this held even when individual agents in the delegation
chain were compromised---the trust bound is a \textit{structural}
guarantee independent of agent behavior.
\end{block}

\begin{block}{Finding 3: Integrity Improvement Scales Across
Architectures}
\protect\phantomsection\label{finding-3-integrity-improvement-scales-across-architectures}
\textbf{Table: Cross-architecture detection summary (Full CIF).}
\{\#tab:cross-arch-summary\}

{\def\LTcaptype{none} % do not increment counter
\begin{longtable}[]{@{}
  >{\raggedright\arraybackslash}p{(\linewidth - 8\tabcolsep) * \real{0.2000}}
  >{\raggedright\arraybackslash}p{(\linewidth - 8\tabcolsep) * \real{0.2000}}
  >{\raggedright\arraybackslash}p{(\linewidth - 8\tabcolsep) * \real{0.2000}}
  >{\raggedright\arraybackslash}p{(\linewidth - 8\tabcolsep) * \real{0.2000}}
  >{\raggedright\arraybackslash}p{(\linewidth - 8\tabcolsep) * \real{0.2000}}@{}}
\toprule\noalign{}
\begin{minipage}[b]{\linewidth}\raggedright
Architecture
\end{minipage} & \begin{minipage}[b]{\linewidth}\raggedright
Overall TPR
\end{minipage} & \begin{minipage}[b]{\linewidth}\raggedright
Strongest Category
\end{minipage} & \begin{minipage}[b]{\linewidth}\raggedright
Weakest Category
\end{minipage} & \begin{minipage}[b]{\linewidth}\raggedright
Latency Overhead
\end{minipage} \\
\midrule\noalign{}
\endhead
Claude Code & 0.94 & Direct injection (0.97) & Coordination (0.88) &
+16\% (p50) \\
AutoGPT & 0.94 & Direct injection (0.96) & Coordination (0.85) & +21\%
(p50) \\
CrewAI & 0.96 & Direct injection (0.97) & Coordination (0.91) & +18\%
(p50)\(^\dagger\) \\
LangGraph & 0.98 & Direct injection (0.98) & Coordination (0.92) & +15\%
(p50)\(^\dagger\) \\
MetaGPT & 0.95 & Direct injection (0.95) & Coordination (0.89) & +19\%
(p50)\(^\dagger\) \\
Camel & 0.92 & Coordination (0.93) & Nested injection (0.89) & +22\%
(p50)\(^\dagger\) \\
\bottomrule\noalign{}
\end{longtable}
}

\(^\dagger\)\textit{Estimated from architecture-specific adapter overhead characteristics.}

CIF improved detection rates across all six architectures, with
particularly strong results for systems with deeper delegation
hierarchies (AutoGPT) where the trust calculus provides the greatest
benefit. Camel's debate architecture is the only system where
coordination attack detection exceeds injection detection---its
adversarial design creates inherent resilience through mutual challenge,
while the peer-to-peer topology exposes it to lateral injection.

The peer-to-peer architectures (Camel) showed the largest relative
improvement, consistent with Part 1's analysis that equal-trust
topologies are most vulnerable to lateral movement attacks
(\cref{tab:architecture-insights}).
\end{block}

\begin{block}{Finding 4: Performance Overhead Is Acceptable for Security
Contexts}
\protect\phantomsection\label{finding-4-performance-overhead-is-acceptable-for-security-contexts}
Full CIF deployment introduces latency overhead in the 20-25\% range
with memory requirements scaling with agent count. For security-critical
deployments, this overhead is acceptable given the integrity improvement
achieved.

The overhead is dominated by the cognitive firewall (input
classification) and Byzantine consensus (coordination). For environments
where consensus is unnecessary, lighter configurations achieve
comparable detection with lower overhead (Table 3 (Risk-Based
Configuration, Part 1)).
\end{block}

\begin{block}{Finding 5: Attack-Type Specific Vulnerabilities Remain}
\protect\phantomsection\label{finding-5-attack-type-specific-vulnerabilities-remain}
Despite strong overall performance, specific attack types remain
challenging:

\begin{itemize}
\item **Semantic equivalent attacks**: Rephrased injections that preserve meaning evade pattern-matching
\item **Progressive drift**: Sub-threshold changes accumulate below detection windows
\item **Orchestrator compromise**: Outside our threat model (our honest orchestrator assumption (Part 1, Section 2))
\end{itemize}

These gaps define the frontier for future defense research.
\end{block}
\end{block}

\begin{block}{Structural Guarantees Beyond Detection Rates}
\protect\phantomsection\label{structural-guarantees-beyond-detection-rates}
The empirical results validate that CIF's formal mechanisms translate to
practical protection. The key insight is not the specific detection
rates achieved---which reflect current attack sophistication and will
degrade as adversaries adapt---but rather the \textit{structural}
properties:

\begin{enumerate}
\item Trust cannot be amplified through delegation (Part 1, Theorem 2)
\item Defenses compose predictably (Part 1, Theorems 3.1 and 3.2)
\item Information-theoretic bounds constrain the stealth-impact tradeoff (Part 1, Theorem 4)
\end{enumerate}

These properties hold independent of specific detection thresholds and
provide the foundation for long-term security assurance.

For detailed statistical analysis including significance testing,
confidence intervals, ablation studies, and scalability benchmarks, see
the Extended Results (\cref{sec:extended-results}).
\end{block}
\end{frame}

\end{document}
